% chapter03.tex — Functions
% Book pp. 39–55 (appendix included). Slice pp. 1–3 are Part II + Dedekind
% (those live in main.tex). No figures.

\spivakchapter{3}{Functions}

Undoubtedly the most important concept in all of mathematics is that of a
function---in almost every branch of modern mathematics functions turn out to
be the central objects of investigation. It will therefore probably not surprise
you to learn that the concept of a function is one of great generality. Perhaps
it will be a relief to learn that, for the present, we will be able to restrict our
attention to functions of a very special kind; even this small class of functions
will exhibit sufficient variety to engage our attention for quite some time. We
will not even begin with a proper definition. For the moment a provisional
definition will enable us to discuss functions at length, and will illustrate the
intuitive notion of functions, as understood by mathematicians. Later, we will
consider and discuss the advantages of the modern mathematical definition.
Let us therefore begin with the following:

\begin{spivakdefn}[PROVISIONAL DEFINITION]
A function is a rule which assigns, to each of certain real numbers, some other
real number.
\end{spivakdefn}

The following examples of functions are meant to illustrate and amplify this
definition, which, admittedly, requires some such clarification.

\textit{Example 1}\quad The rule which assigns to each number the square of that
number.

\textit{Example 2}\quad The rule which assigns to each number $y$ the number
\[
\frac{y^3 + 3y + 5}{y^2 + 1}.
\]

\textit{Example 3}\quad The rule which assigns to each number $c \neq 1$, $-1$ the
number
\[
\frac{c^3 + 3c + 5}{c^2 - 1}.
\]

\textit{Example 4}\quad The rule which assigns to each number $x$ satisfying
$-17 \leq x \leq \pi/3$ the number $x^2$.

\textit{Example 5}\quad The rule which assigns to each number $a$ the number $0$
if $a$ is irrational, and the number $1$ if $a$ is rational.

\textit{Example 6}\quad The rule which assigns
\begin{align*}
&\text{to $2$ the number $5$,}\\
&\text{to $17$ the number $\dfrac{36}{\pi}$,}\\
&\text{to $\dfrac{\pi^2}{17}$ the number $28$,}\\
&\text{to $\dfrac{36}{\pi}$ the number $28$,}
\end{align*}
and to any $y \neq 2$, $17$, $\pi^2/17$, or $36/\pi$, the number $16$ if $y$ is of
the form $a + b\sqrt{2}$ for $a$, $b$ in $\mathbb{Q}$.

\textit{Example 7}\quad The rule which assigns to each number $t$ the number
$t^3 + x$. (This rule depends, of course, on what the number $x$ is, so we are
really describing infinitely many different functions, one for each number $x$.)

\textit{Example 8}\quad The rule which assigns to each number $z$ the number of
$7$'s in the decimal expansion of $z$, if this number is finite, and $-\pi$ if there
are infinitely many $7$'s in the decimal expansion of $z$.

One thing should be abundantly clear from these examples---a function is
\emph{any} rule that assigns numbers to certain other numbers, not just a rule
which can be expressed by an algebraic formula, or even by one uniform
condition which applies to every number; nor is it necessarily a rule which you,
or anybody else, can actually apply in practice (no one knows, for example,
what rule $8$ associates to $\pi$). Moreover, the rule may neglect some numbers
and it may not even be clear to which numbers the function applies (try to
determine, for example, whether the function in Example~$6$ applies to $\pi$).
The set of numbers to which a function \emph{does} apply is called the
\textbf{domain} of the function.

Before saying anything else about functions we badly need some notation.
Since throughout this book we shall frequently be talking about functions
(indeed we shall hardly ever talk about anything else) we need a convenient
way of naming functions, and of referring to functions in general. The
standard practice is to denote a function by a letter. For obvious reasons the
letter ``$f$'' is a favorite, thereby making ``$g$'' and ``$h$'' other obvious
candidates, but any letter (or any reasonable symbol, for that matter) will do,
not excluding ``$x$'' and ``$y$'', although these letters are usually reserved for
indicating numbers. If $f$ is a function, then the number which $f$ associates
to a number $x$ is denoted by $f(x)$---this symbol is read ``$f$ of $x$'' and is
often called the \textbf{value of $\boldsymbol{f}$ at $\boldsymbol{x}$}. Naturally, if we denote a
function by $x$, some other letter must be chosen to denote the number (a
perfectly legitimate, though perverse, choice would be ``$f$,'' leading to the
symbol $x(f)$). Note that the symbol $f(x)$ makes sense only for $x$ in the
domain of $f$; for other $x$ the symbol $f(x)$ is not defined.

If the functions defined in Examples~$1$--$8$ are denoted by $f$, $g$, $h$, $r$,
$s$, $\theta$, $\alpha_x$, and $y$, then we can rewrite their definitions as follows:
\begin{steps}
\item $f(x) = x^2$ \quad for all $x$.
\item $g(y) = \dfrac{y^3 + 3y + 5}{y^2 + 1}$ \quad for all $y$.
\item $h(c) = \dfrac{c^3 + 3c + 5}{c^2 - 1}$ \quad for all $c \neq 1$, $-1$.
\item $r(x) = x^2$ \quad for all $x$ such that $-17 \leq x \leq \pi/3$.
\item $\displaystyle s(x) =
\begin{cases}
0, & x\text{ irrational}\\
1, & x\text{ rational.}
\end{cases}$
\item $\displaystyle \theta(x) =
\begin{cases}
5, & x = 2\\[4pt]
\dfrac{36}{\pi}, & x = 17\\[8pt]
28, & x = \dfrac{\pi^2}{17}\\[8pt]
28, & x = \dfrac{36}{\pi}\\[8pt]
16, &
\begin{aligned}[t]
&x \neq 2,\ 17,\ \dfrac{\pi^2}{17},\ \text{or }\dfrac{36}{\pi},\\
&\text{and }x = a + b\sqrt{2}\text{ for }a,\ b\text{ in }\mathbb{Q}.
\end{aligned}
\end{cases}$
\item $\alpha_x(t) = t^3 + x$ \quad for all numbers $t$.
\item $\displaystyle y(x) =
\begin{cases}
n, & \text{exactly $n$ $7$'s appear in the decimal expansion of $x$}\\
-\pi, & \text{infinitely many $7$'s appear in the decimal expansion of $x$.}
\end{cases}$
\end{steps}

These definitions illustrate the common procedure adopted for defining a
function $f$---indicating what $f(x)$ is for every number $x$ in the domain of
$f$. (Notice that this is exactly the same as indicating $f(a)$ for every number
$a$, or $f(b)$ for every number $b$, etc.) In practice, certain abbreviations are
tolerated. Definition~(1) could be written simple
\[
(1)\qquad f(x) = x^2
\]
the qualifying phrase ``for all $x$'' being understood. Of course, for
definition~(4) the only possible abbreviation is
\[
(4)\qquad r(x) = x^2, \qquad -17 \leq x \leq \pi/3.
\]
It is usually understood that a definition such as
\[
k(x) = \frac{1}{x} + \frac{1}{x-1}, \qquad x \neq 0,\ 1
\]
can be shortened to
\[
k(x) = \frac{1}{x} + \frac{1}{x-1};
\]
in other words, \emph{unless the domain is explicitly restricted further, it is
understood to consist of all numbers for which the definition makes any sense
at all.}

You should have little difficulty checking the following assertions about the
functions defined above:
\begin{align*}
f(x+1) &= f(x) + 2x + 1;\\
g(x) &= h(x)\text{ if }x^3 + 3x + 5 = 0;\\
r(x+1) &= r(x) + 2x + 1\text{ if }{-17} \leq x \leq \tfrac{\pi}{3} - 1;\\
s(x+y) &= s(x)\text{ if $y$ is rational;}\\
\theta\Bigl(\frac{\pi^2}{17}\Bigr) &= \theta\Bigl(\frac{36}{\pi}\Bigr);\\
\alpha_x(x) &= x \cdot \bigl[f(x) + 1\bigr];\\
y\Bigl(\frac{1}{3}\Bigr) &= 0,\quad
y\Bigl(\frac{7}{9}\Bigr) = -\pi.
\end{align*}

If the expression $f(s(a))$ looks unreasonable to you, then you are forgetting
that $s(a)$ is a number like any other number, so that $f(s(a))$ makes sense.
As a matter of fact, $f(s(a)) = s(a)$ for all $a$. Why? Even more complicated
expressions than $f(s(a))$ are, after a first exposure, no more difficult to
unravel. The expression
\[
f\bigl(r\bigl(s\bigl(\theta\bigl(\alpha_3\bigl(y(\tfrac13)\bigr)\bigr)\bigr)\bigr)\bigr),
\]
formidable as it appears, may be evaluated quite easily with a little patience:
\begin{align*}
f\bigl(r\bigl(s\bigl(\theta\bigl(\alpha_3\bigl(y(\tfrac13)\bigr)\bigr)\bigr)\bigr)\bigr)
&= f\bigl(r\bigl(s\bigl(\theta\bigl(\alpha_3(0)\bigr)\bigr)\bigr)\bigr)\\
&= f\bigl(r\bigl(s\bigl(\theta(3)\bigr)\bigr)\bigr)\\
&= f\bigl(r\bigl(s(16)\bigr)\bigr)\\
&= f\bigl(r(1)\bigr)\\
&= f(1)\\
&= 1.
\end{align*}

The first few problems at the end of this chapter give further practice
manipulating this symbolism.

The function defined in~(1) is a rather special example of an extremely
important class of functions, the polynomial functions. A function $f$ is a
\textbf{polynomial function} if there are real numbers $a_0, \dots, a_n$ such
that
\[
f(x) = a_n x^n + a_{n-1}x^{n-1} + \cdots + a_2 x^2 + a_1 x + a_0,
\qquad\text{for all $x$}
\]
(when $f(x)$ is written in this form it is usually tacitly assumed that
$a_n \neq 0$). The highest power of $x$ with a nonzero coefficient is called
the \textbf{degree} of $f$; for example, the polynomial function $f$ defined by
$f(x) = 5x^6 + 137x^4 - \pi$ has degree~$6$.

The functions defined in~(2) and~(3) belong to a somewhat larger class of
functions, the \textbf{rational functions}; these are the functions of the form
$p/q$ where $p$ and $q$ are polynomial functions (and $q$ is not the function
which is always~$0$). The rational functions are themselves quite special
examples of an even larger class of functions, very thoroughly studied in
calculus, which are simpler than many of the functions first mentioned in this
chapter. The following are examples of this kind of function:
\begin{steps}[start=9]
\item $f(x) = \dfrac{x + x^2 + x\sin^2 x}{x\sin x + x\sin^2 x}$
\item $f(x) = \sin(x^2)$.
\item $f(x) = \sin(\sin(x^2))$.
\item $\displaystyle
f(x) = \sin^2\bigl(\sin(\sin^2(x\sin^2 x^2))\bigr)
\cdot \sin\Biggl(\frac{x + \sin(x\sin x)}{x + \sin x}\Biggr).$
\end{steps}

By what criterion, you may feel impelled to ask, can such functions,
especially a monstrosity like~(12), be considered simple? The answer is that
they can be built up from a few simple functions using a few simple means of
combining functions. In order to construct the functions~(9)--(12) we need to
start with the ``identity function'' $I$, for which $I(x) = x$, and the ``sine
function'' $\sin$, whose value $\sin(x)$ at $x$ is often written simple
$\sin x$. The following are some of the important ways in which functions may
be combined to produce new functions.

If $f$ and $g$ are any two functions, we can define a new function $f+g$,
called the \textbf{sum} of $f$ and $g$, by the equation
\[
(f+g)(x) = f(x) + g(x).
\]
Note that according to the conventions we have adopted, the domain of $f+g$
consists of all $x$ for which ``$f(x)+g(x)$'' makes sense, i.e., the set of all
$x$ in both $\dom f$ and $\dom g$. If $A$ and $B$ are any two sets, then
$A \cap B$ (read ``$A$ intersect $B$'' or ``the intersection of $A$ and $B$'')
denotes the set of $x$ in both $A$ and $B$; this notation allows us to write
$\dom(f+g) = \dom f \cap \dom g$.

In a similar vein, we define the \textbf{product} $f \cdot g$ and the
\textbf{quotient} $\dfrac{f}{g}$ (or $f/g$) of $f$ and $g$ by
\[
(f \cdot g)(x) = f(x) \cdot g(x)
\]
and
\[
\Biggl(\frac{f}{g}\Biggr)(x) = \frac{f(x)}{g(x)}.
\]
Moreover, if $g$ is a function and $c$ is a number, we define a new function
$c \cdot g$ by
\[
(c \cdot g)(x) = c \cdot g(x).
\]
This becomes a special case of the notation $f \cdot g$ if we agree that the
symbol $c$ should also represent the function $f$ defined by $f(x) = c$; such
a function, which has the same value for all numbers $x$, is called a
\textbf{constant function}.

The domain of $f \cdot g$ is $\dom f \cap \dom g$, and the domain of
$c \cdot g$ is simply the domain of $g$. On the other hand, the domain of
$f/g$ is rather complicated---it may be written
$\dom f \cap \dom g \cap \{x : g(x) \neq 0\}$, the symbol
$\{x : g(x) \neq 0\}$ denoting the set of numbers $x$ such that
$g(x) \neq 0$. In general, $\{x : \ldots\}$ denotes the set of all $x$ such
that ``$\ldots$'' is true. Thus $\{x : x^3 + 3 < 11\}$ denotes the set of all
numbers $x$ such that $x^3 < 8$, and consequently
$\{x : x^3 + 3 < 11\} = \{x : x < 2\}$. Either of these symbols could just as
well have been written using $y$ everywhere instead of $x$. Variations of this
notation are common, but hardly require any discussion. Any one can guess
that $\{x > 0 : x^3 < 8\}$ denotes the set of positive numbers whose cube is
less than~$8$; it could be expressed more formally as
$\{x : x > 0\text{ and }x^3 < 8\}$. Incidentally, this set is equal to the set
$\{x : 0 < x < 2\}$. One variation is slightly less transparent, but very
standard. The set $\{1, 3, 2, 4\}$, for example, contains just the four numbers
$1$, $2$, $3$, and~$4$; it can also be denoted by
$\{x : x = 1\text{ or }x = 3\text{ or }x = 2\text{ or }x = 4\}$.

Certain facts about the sum, product, and quotient of functions are obvious
consequences of facts about sums, products, and quotients of numbers. For
example, it is very easy to prove that
\[
(f+g)+h = f+(g+h).
\]
The proof is characteristic of almost every proof which demonstrates that two
functions are equal---the two functions must be shown to have the same domain,
and the same value at any number in the domain. For example, to prove that
$(f+g)+h = f+(g+h)$, note that unraveling the definition of the two sides
gives
\begin{align*}
[(f+g)+h](x)
&= (f+g)(x) + h(x)\\
&= \bigl[f(x) + g(x)\bigr] + h(x)
\end{align*}
and
\begin{align*}
[f+(g+h)](x)
&= f(x) + (g+h)(x)\\
&= f(x) + \bigl[g(x) + h(x)\bigr],
\end{align*}
and the equality of $[f(x)+g(x)]+h(x)$ and $f(x)+[g(x)+h(x)]$ is a fact about
numbers. In this proof the equality of the two domains was not explicitly
mentioned because this is obvious, as soon as we begin to write down these
equations; the domain of $(f+g)+h$ and of $f+(g+h)$ is clearly
$\dom f \cap \dom g \cap \dom h$. We naturally write $f+g+h$ for
$(f+g)+h = f+(g+h)$, precisely as we did for numbers.

It is just as easy to prove that $(f \cdot g) \cdot h = f \cdot (g \cdot h)$,
and this function is denoted by $f \cdot g \cdot h$. The equations
$f+g = g+f$ and $f \cdot g = g \cdot f$ should also present no difficulty.

Using the operations $+$, $\cdot$, $/$ we can now express the function $f$
defined in~(9) by
\[
f = \frac{I + I \cdot I + I \cdot \sin \cdot \sin}
{I \cdot \sin + I \cdot \sin \cdot \sin}.
\]
It should be clear, however, that we cannot express function~(10) this way. We
require yet another way of combining functions. This combination, the
composition of two functions, is by far the most important.

If $f$ and $g$ are any two functions, we define a new function $f \circ g$,
the \textbf{composition} of $f$ and $g$, by
\[
(f \circ g)(x) = f(g(x));
\]
the domain of $f \circ g$ is $\{x : x\text{ is in }\dom g\text{ and }g(x)
\text{ is in }\dom f\}$. The symbol ``$f \circ g$'' is often read ``$f$ circle
$g$.'' Compared to the phrase ``the composition of $f$ and $g$'' this has the
advantage of brevity, of course, but there is another advantage of far greater
import: there is much less chance of confusing $f \circ g$ with $g \circ f$,
and these must \emph{not} be confused, since they are not usually equal; in
fact, almost any $f$ and $g$ chosen at random will illustrate this point (try
$f = I \cdot I$ and $g = \sin$, for example). Lest you become too apprehensive
about the operation of composition, let us hasten to point out that
composition \emph{is} associative:
\[
(f \circ g) \circ h = f \circ (g \circ h)
\]
(and the proof is a triviality); this function is denoted by $f \circ g \circ h$.
We can now write the functions~(10), (11), (12) as
\begin{align*}
(10)\quad f &= \sin \circ (I \cdot I),\\
(11)\quad f &= \sin \circ \sin \circ (I \cdot I),\\
(12)\quad f &= (\sin \cdot \sin) \circ \sin \circ (\sin \cdot \sin)
\circ \bigl(I \cdot [(\sin \cdot \sin) \circ (I \cdot I)]\bigr)\\
&\qquad\cdot
\sin \circ \Biggl(\frac{I + \sin \circ (I \cdot \sin)}{I + \sin}\Biggr).
\end{align*}

One fact has probably already become clear. Although this method of writing
functions reveals their ``structure'' very clearly, it is hardly short or
convenient. The shortest name for the function $f$ such that
$f(x) = \sin(x^2)$ for all $x$ unfortunately seems to be ``the function $f$
such that $f(x) = \sin(x^2)$ for all $x$.'' The need for abbreviating this
clumsy description has been clear for two hundred years, but no reasonable
abbreviation has received universal acclaim. At present the strongest
contender for this honor is something like
\[
x \to \sin(x^2)
\]
(read ``$x$ goes to $\sin(x^2)$'' or just ``$x$ arrow $\sin(x^2)$''), but it is
hardly popular among writers of calculus textbooks. In this book we will
tolerate a certain amount of ellipsis, and speak of ``the function
$f(x) = \sin(x^2)$.'' Even more popular is the quite drastic abbreviation:
``the function $\sin(x^2)$.'' For the sake of precision we will never use this
description, which, strictly speaking, confuses a number and a function, but
it is so convenient that you will probably end up adopting it for personal
use. As with any convention, utility is the motivating factor, and this
criterion is reasonable so long as the slight logical deficiencies cause no
confusion. On occasion, confusion \emph{will} arise unless a more precise
description is used. For example, ``the function $x + t^3$'' is an ambiguous
phrase; it could mean either
\[
x \to x + t^3,\text{ i.e., the function $f$ such that $f(x) = x + t^3$
for all $x$}
\]
or
\[
t \to x + t^3,\text{ i.e., the function $f$ such that $f(t) = x + t^3$
for all $t$}.
\]
As we shall see, however, for many important concepts associated with
functions, calculus has a notation which contains the ``$x \to$'' built in.

By now we have made a sufficiently extensive investigation of functions to
warrant reconsidering our definition. We have defined a function as a
``rule,'' but it is hardly clear what this means. If we ask ``What happens if
you break this rule?'' it is not easy to say whether this question is merely
facetious or actually profound.

A more substantial objection to the use of the word ``rule'' is that
\[
f(x) = x^2
\]
and
\[
f(x) = x^2 + 3x + 3 - 3(x+1)
\]
are certainly \emph{different} rules, if by a rule we mean the actual
instructions given for determining $f(x)$; nevertheless, we want
\[
f(x) = x^2
\]
and
\[
f(x) = x^2 + 3x + 3 - c(x+1)
\]
to define the same function. For this reason, a function is sometimes defined
as an ``association'' between numbers; unfortunately the word ``association''
escapes the objections raised against ``rule'' only because it is even more
vague.

There is, of course, a satisfactory way of defining functions, or we should
never have gone to the trouble of criticizing our original definition. But a
satisfactory definition can never be constructed by finding synonyms for
English words which are troublesome. The definition which mathematicians have
finally accepted for ``function'' is a beautiful example of the means by which
intuitive ideas have been incorporated into rigorous mathematics. The correct
question to ask about a function is not ``What is a rule?'' or ``What is an
association?'' but ``What does one have to know about a function in order to
know all about it?'' The answer to the last question is easy---for each number
$x$ one needs to know the number $f(x)$; we can imagine a table which would
display all the information one could desire about the function
$f(x) = x^2$:
\[
\begin{array}{cc}
x & f(x)\\
\hline
1 & 1\\
-1 & 1\\
2 & 4\\
-2 & 4\\
\sqrt{2} & 2\\
-\sqrt{2} & 2\\
\pi & \pi^2\\
-\pi & \pi^2
\end{array}
\]
It is not even necessary to arrange the numbers in a table (which would
actually be impossible if we wanted to list all of them). Instead of a two
column array we can consider various pairs of numbers
\[
(1,1),\ (-1,1),\ (2,4),\ (-2,4),\ (\pi,\pi^2),\ (\sqrt{2},2),\dots
\]
simply collected together into a set.\footnote{The pairs occurring here are
often called ``ordered pairs,'' to emphasize that, for example, $(2,4)$ is not
the same pair as $(4,2)$. It is only fair to warn that we are going to define
functions in terms of ordered pairs, another undefined term. Ordered pairs
\emph{can} be defined, however, and an appendix to this chapter has been
provided for skeptics.} To find $f(1)$ we simply take the second number of
the pair whose first member is~$1$; to find $f(\pi)$ we take the second number
of the pair whose first member is~$\pi$. We seem to be saying that a function
might as well be defined as a collection of pairs of numbers. For example, if
we were given the following collection (which contains just~$5$ pairs):
\[
f = \bigl\{(1,7),\ (3,7),\ (5,3),\ (4,8),\ (8,4)\bigr\},
\]
then $f(1)=7$, $f(3)=7$, $f(5)=3$, $f(4)=8$, $f(8)=4$ and $1$, $3$, $4$, $5$,
$8$ are the only numbers in the domain of $f$. If we consider the collection
\[
f = \bigl\{(1,7),\ (3,7),\ (2,5),\ (1,8),\ (8,4)\bigr\},
\]
then $f(3)=7$, $f(2)=5$, $f(8)=4$; but it is impossible to decide whether
$f(1)=7$ or $f(1)=8$. In other words, a function cannot be defined to be any
old collection of pairs of numbers; we must rule out the possibility which
arose in this case. We are therefore led to the following definition.

\begin{spivakdefn}
A \textbf{function} is a collection of pairs of numbers with the following
property: if $(a,b)$ and $(a,c)$ are both in the collection, then $b=c$; in
other words, the collection must not contain two different pairs with the same
first element.
\end{spivakdefn}

This is our first full-fledged definition, and illustrates the format we shall
always use to define significant new concepts. These definitions are so
important (at least as important as theorems) that it is essential to know
when one is actually at hand, and to distinguish them from comments,
motivating remarks, and casual explanations. They will be preceded by the
word \textbf{DEFINITION}, contain the term being defined in boldface letters,
and constitute a paragraph unto themselves.

There is one more definition (actually defining two things at once) which can
now be made rigorously:

\begin{spivakdefn}
If $f$ is a function, the \textbf{domain} of $f$ is the set of all $a$ for
which there is some $b$ such that $(a,b)$ is in $f$. If $a$ is in the domain
of $f$, it follows from the definition of a function that there is, in fact, a
\emph{unique} number $b$ such that $(a,b)$ is in $f$. This unique $b$ is
denoted by $\boldsymbol{f(a)}$.
\end{spivakdefn}

With this definition we have reached our goal: the important thing about a
function $f$ is that a number $f(x)$ is determined for each number $x$ in its
domain. You may feel that we have also reached the point where an intuitive
definition has been replaced by an abstraction with which the mind can hardly
grapple. Two consolations may be offered. First, although a function has been
defined as a collection of pairs, there is nothing to stop you from
\emph{thinking} of a function as a rule. Second, neither the intuitive nor the
formal definition indicates the best way of thinking about functions. The best
way is to draw pictures; but this requires a chapter all by itself.

\subsection*{PROBLEMS}

\pnum{1.} Let $f(x) = 1/(1+x)$. What is
\begin{romans}
\item $f(f(x))$ (for which $x$ does this make sense?).
\item $\displaystyle f\Biggl(\frac{1}{x}\Biggr)$.
\item $f(cx)$.
\item $f(x+y)$.
\item $f(x)+f(y)$.
\item For which numbers $c$ is there a number $x$ such that $f(cx)=f(x)$.
Hint: There are a lot more than you might think at first glance.
\item For which numbers $c$ is it true that $f(cx)=f(x)$ for two different
numbers $x$?
\end{romans}

\pnum{2.} Let $g(x)=x^2$, and let
\[
h(x) =
\begin{cases}
0, & x\text{ rational}\\
1, & x\text{ irrational.}
\end{cases}
\]
\begin{romans}
\item For which $y$ is $h(y) \leq y$?
\item For which $y$ is $h(y) \leq g(y)$?
\item What is $g(h(z)) - h(z)$?
\item For which $w$ is $g(w) \leq w$?
\item For which $\varepsilon$ is $g(g(\varepsilon)) = g(\varepsilon)$?
\end{romans}

\pnum{3.} Find the domain of the functions defined by the following formulas.
\begin{romans}
\item $f(x) = \sqrt{1-x^2}$.
\item $f(x) = \sqrt{1 - \sqrt{1-x^2}}$.
\item $\displaystyle f(x) = \frac{1}{x-1} + \frac{1}{x-2}$.
\item $f(x) = \sqrt{1-x^2} + \sqrt{x^2-1}$.
\item $f(x) = \sqrt{1-x} + \sqrt{x-2}$.
\end{romans}

\pnum{4.} Let $S(x)=x^2$, let $P(x)=2^x$, and let $s(x)=\sin x$. Find each of
the following. In each case you answer should be a \emph{number}.
\begin{romans}
\item $(S \circ P)(y)$.
\item $(S \circ s)(y)$.
\item $(S \circ P \circ s)(t) + (s \circ P)(t)$.
\item $s(t^3)$.
\end{romans}

\pnum{5.} Express each of the following functions in terms of $S$, $P$, $s$,
using only $+$, $\cdot$, and $\circ$ (for example, the answer to~(i) is
$P \circ s$). In each case your answer should be a \emph{function}.
\begin{romans}
\item $f(x) = 2^{\sin x}$.
\item $f(x) = \sin 2^x$.
\item $f(x) = \sin x^2$.
\item $f(x) = \sin^2 x$ (remember that $\sin^2 x$ is an abbreviation for
$(\sin x)^2$).
\item $f(t) = 2^{2^t}$. (Note: $a^{b^c}$ \emph{always} means $a^{(b^c)}$; this
convention is adopted because $(a^b)^c$ can be written more simply as
$a^{bc}$.)
\item $f(u) = \sin(2^u + 2^{u^2})$.
\item $f(y) = \sin\bigl(\sin\bigl(\sin(2^{2^{\sin y}})\bigr)\bigr)$.
\item $f(a) = 2^{\sin^2 a} + \sin(a^2) + 2^{\sin(a^2 + \sin a)}$.
\end{romans}

Polynomial functions, because they are simple, yet flexible, occupy a favored
role in most investigations of functions. The following two problems
illustrate their flexibility, and guide you through a derivation of their most
important elementary properties.

\pnum{6.}
\begin{letters}
\item If $x_1, \dots, x_n$ are distinct numbers, find a polynomial function
$f_i$ of degree $n-1$ which is $1$ at $x_i$ and $0$ at $x_j$ for $j \neq i$.
Hint: the product of all $(x-x_j)$ for $j \neq i$, is $0$ at $x_j$ if
$j \neq i$. (This product is usually denoted by
\[
\prod_{\substack{j=1\\ j\neq i}}^n (x-x_j),
\]
the symbol $\Pi$ (capital pi) playing the same role for products that $\Sigma$
plays for sums.)
\item Now find a polynomial function $f$ of degree $n-1$ such that
$f(x_i)=a_i$, where $a_1, \dots, a_n$ are given numbers. (You should use the
functions $f_i$ from part~(a). The formula you will obtain is called the
``Lagrange interpolation formula.'')
\end{letters}

\pnum{7.}
\begin{letters}
\item Prove that for any polynomial function $f$, and any number $a$, there is
a polynomial function $g$, and a number $b$, such that
$f(x)=(x-a)g(x)+b$ for all $x$. (The idea is simply to divide $(x-a)$ into
$f(x)$ by long division, until a constant remainder is left. For example, the
calculation
\begin{center}
\begin{tabular}{@{}r@{\,}c@{}r@{}r@{}r@{}r@{}}
 && $x^2$ & ${}+x$ && ${}-2$ \\
\cline{3-6}
$x-1$ & $\big)$& $x^3$ && ${}-3x$ & ${}+1$ \\
 && $x^3$ & ${}-x^2$ \\
\cline{3-4}
 &&& $x^2$ & ${}-3x$ \\
 &&& $x^2$ & ${}-x$ \\
\cline{4-5}
 &&&& ${}-2x$ & ${}+1$ \\
 &&&& ${}-2x$ & ${}+2$ \\
\cline{5-6}
 &&&&& ${}-1$
\end{tabular}
\end{center}
shows that $x^3 - 3x + 1 = (x-1)(x^2 + x - 2) - 1$. A formal proof is possible
by induction on the degree of $f$.)
\item Prove that if $f(a)=0$, then $f(x)=(x-a)g(x)$ for some polynomial
function $g$. (The converse is obvious.)
\item Prove that if $f$ is a polynomial function of degree $n$, then $f$ has
at most $n$ roots, i.e., there are at most $n$ numbers $a$ with $f(a)=0$.
\item Show that for each $n$ there is a polynomial function of degree $n$ with
$n$ roots. If $n$ is even find a polynomial function of degree $n$ with no
roots, and if $n$ is odd find one with only one root.
\end{letters}

\pnum{8.} For which numbers $a$, $b$, $c$, and $d$ will the function
\[
f(x) = \frac{ax+b}{cx+d}
\]
satisfy $f(f(x))=x$ for all $x$ (for which this equation makes sense)?

\pnum{9.}
\begin{letters}
\item If $A$ is any set of real numbers, define a function $C_A$ as follows:
\[
C_A(x) =
\begin{cases}
1, & x\text{ in }A\\
0, & x\text{ not in }A.
\end{cases}
\]
Find expressions for $C_{A\cap B}$ and $C_{A\cup B}$ and $C_{\mathbf{R}-A}$,
in terms of $C_A$ and $C_B$. (The symbol $A \cap B$ was defined in this
chapter, but the other two may be new to you. They can be defined as follows:
\begin{align*}
A \cup B &= \{x : x\text{ is in }A\text{ or }x\text{ is in }B\},\\
\mathbf{R} - A &= \{x : x\text{ is in }\mathbf{R}\text{ but }x\text{ is not
in }A\}.)
\end{align*}
\item Suppose $f$ is a function such that $f(x)=0$ or $1$ for each $x$. Prove
that there is a set $A$ such that $f = C_A$.
\item Show that $f = f^2$ if and only if $f = C_A$ for some set $A$.
\end{letters}

\pnum{10.}
\begin{letters}
\item For which functions $f$ is there a function $g$ such that $f = g^2$?
Hint: You can certainly answer this question if ``function'' is replaced by
``number.''
\item For which functions $f$ is there a function $g$ such that $f = 1/g$?
\item[*(c)] For which functions $b$ and $c$ can we find a function $x$ such
that
\[
(x(t))^2 + b(t)x(t) + c(t) = 0
\]
for all numbers $t$?
\item[*(d)] What conditions must the functions $a$ and $b$ satisfy if there
is to be a function $x$ such that
\[
a(t)x(t) + b(t) = 0
\]
for all numbers $t$? How many such functions $x$ will there be?
\end{letters}

\pnum{11.}
\begin{letters}
\item Suppose that $H$ is a function and $y$ is a number such that
$H(H(y))=y$. What is
\[
\underbrace{H(H(H(\cdots(H(y)\cdots)}_{80\text{ times}}?
\]
\item Same question if $80$ is replaced by $81$.
\item Same question if $H(H(y))=H(y)$.
\item[*(d)] Find a function $H$ such that $H(H(x))=H(x)$ for all numbers $x$,
and such that $H(1)=36$, $H(2)=\pi/3$, $H(13)=47$, $H(36)=36$,
$H(\pi/3)=\pi/3$, $H(47)=47$. (Don't try to ``solve'' for $H(x)$; there are
many functions $H$ with $H(H(x))=H(x)$. The extra conditions on $H$ are
supposed to suggest a way of finding a suitable $H$.)
\item[*(e)] Find a function $H$ such that $H(H(x))=H(x)$ for all $x$, and such
that $H(1)=7$, $H(17)=18$.
\end{letters}

\pnum{12.} A function $f$ is \textbf{even} if $f(x)=f(-x)$ and \textbf{odd} if
$f(x)=-f(-x)$. For example, $f$ is even if $f(x)=x^2$ or $f(x)=\abs{x}$ or
$f(x)=\cos x$, while $f$ is odd if $f(x)=x$ or $f(x)=\sin x$.
\begin{letters}
\item Determine whether $f+g$ is even, odd, or not necessarily either, in the
four cases obtained by choosing $f$ even or odd, and $g$ even or odd. (Your
answers can most conveniently be displayed in a $2 \times 2$ table.)
\item Do the same for $f \cdot g$.
\item Do the same for $f \circ g$.
\item Prove that every even function $f$ can be written $f(x)=g(\abs{x})$, for
infinitely many functions $g$.
\end{letters}

\pnum{*13.}
\begin{letters}
\item Prove that any function $f$ with domain $\mathbf{R}$ can be written
$f = E + O$, where $E$ is even and $O$ is odd.
\item Prove that this way of writing $f$ is unique. (If you try to do
part~(b) first, by ``solving'' for $E$ and $O$ you will probably find the
solution to part~(a).)
\end{letters}

\pnum{14.} If $f$ is any function, define a new function $\abs{f}$ by
$\abs{f}(x)=\abs{f(x)}$. If $f$ and $g$ are functions, define two new
functions, $\max(f,g)$ and $\min(f,g)$, by
\begin{align*}
\max(f,g)(x) &= \max(f(x), g(x)),\\
\min(f,g)(x) &= \min(f(x), g(x)).
\end{align*}
Find an expression for $\max(f,g)$ and $\min(f,g)$ in terms of $\abs{\ }$.

\pnum{15.}
\begin{letters}
\item Show that $f = \max(f,0) + \min(f,0)$. This particular way of writing
$f$ is fairly useful; the functions $\max(f,0)$ and $\min(f,0)$ are called the
\textbf{positive} and \textbf{negative parts} of $f$.
\item A function $f$ is called \textbf{nonnegative} if $f(x) \geq 0$ for all
$x$. Prove that any function $f$ can be written $f = g - h$, where $g$ and $h$
are nonnegative, in infinitely many ways. (The ``standard way'' is
$g = \max(f,0)$ and $h = -\min(f,0)$.) Hint: Any \emph{number} can certainly
be written as the difference of two nonnegative \emph{numbers} in infinitely
many ways.
\end{letters}

\pnum{*16.} Suppose $f$ satisfies $f(x+y)=f(x)+f(y)$ for all $x$ and $y$.
\begin{letters}
\item Prove that $f(x_1 + \cdots + x_n) = f(x_1) + \cdots + f(x_n)$.
\item Prove that there is some number $c$ such that $f(x)=cx$ for all
\emph{rational} numbers $x$ (at this point we're not trying to say anything
about $f(x)$ for irrational $x$). Hint: First figure out what $c$ must be. Now
prove that $f(x)=cx$, first when $x$ is a natural number, then when $x$ is an
integer, then when $x$ is the reciprocal of an integer and, finally, for all
rational $x$.
\end{letters}

\pnum{*17.} If $f(x)=0$ for all $x$, then $f$ satisfies
$f(x+y)=f(x)+f(y)$ for all $x$ and $y$, and also $f(x \cdot y)=f(x) \cdot f(y)$
for all $x$ and $y$. Now suppose that $f$ satisfies these two properties, but
that $f(x)$ is not always~$0$. Prove that $f(x)=x$ for all $x$, as follows:
\begin{letters}
\item Prove that $f(1)=1$.
\item Prove that $f(x)=x$ if $x$ is rational.
\item Prove that $f(x)>0$ if $x>0$. (This part is tricky, but if you have been
paying attention to the philosophical remarks accompanying the problems in the
last two chapters, you will know what to do.)
\item Prove that $f(x)>f(y)$ if $x>y$.
\item Prove that $f(x)=x$ for all $x$. Hint: Use the fact that between any two
numbers there is a rational number.
\end{letters}

\pnum{*18.} Precisely what conditions must $f$, $g$, $h$, and $k$ satisfy in
order that $f(x)g(y)=h(x)k(y)$ for all $x$ and $y$?

\pnum{*19.}
\begin{letters}
\item Prove that there do \emph{not} exist functions $f$ and $g$ with either
of the following properties:
\begin{romans}
\item $f(x)+g(y)=xy$ for all $x$ and $y$.
\item $f(x) \cdot g(y)=x+y$ for all $x$ and $y$.
\end{romans}
Hint: Try to get some information about $f$ or $g$ by choosing particular
values of $x$ and $y$.
\item Find functions $f$ and $g$ such that $f(x+y)=g(xy)$ for all $x$ and $y$.
\end{letters}

\pnum{*20.}
\begin{letters}
\item Find a function $f$, other than a constant function, such that
$\abs{f(y)-f(x)} \leq \abs{y-x}$.
\item Suppose that $f(y)-f(x) \leq (y-x)^2$ for all $x$ and $y$. (Why does this
imply that $\abs{f(y)-f(x)} \leq (y-x)^2$?) Prove that $f$ is a constant
function. Hint: Divide the interval from $x$ to $y$ into $n$ equal pieces.
\end{letters}

\pnum{21.} Prove or give a counterexample for each of the following
assertions:
\begin{letters}
\item $f \circ (g+h) = f \circ g + f \circ h$.
\item $(g+h) \circ f = g \circ f + h \circ f$.
\item $\dfrac{1}{f \circ g} = \dfrac{1}{f} \circ g$.
\item $\dfrac{1}{f \circ g} = f \circ \Biggl(\dfrac{1}{g}\Biggr)$.
\end{letters}

\pnum{22.}
\begin{letters}
\item Suppose $g = h \circ f$. Prove that if $f(x)=f(y)$, then $g(x)=g(y)$.
\item Conversely, suppose that $f$ and $g$ are two functions such that
$g(x)=g(y)$ whenever $f(x)=f(y)$. Prove that $g = h \circ f$ for some function
$h$. Hint: Just try to define $h(z)$ when $z$ is of the form $z=f(x)$ (these
are the only $z$ that matter) and use the hypotheses to show that your
definition will not run into trouble.
\end{letters}

\pnum{23.} Suppose that $f \circ g = I$, where $I(x)=x$. Prove that
\begin{letters}
\item if $x \neq y$, then $g(x) \neq g(y)$;
\item every number $b$ can be written $b=f(a)$ for some number $a$.
\end{letters}

\pnum{*24.}
\begin{letters}
\item Suppose $g$ is a function with the property that $g(x) \neq g(y)$ if
$x \neq y$. Prove that there is a function $f$ such that $f \circ g = I$.
\item Suppose that $f$ is a function such that every number $b$ can be written
$b=f(a)$ for some number $a$. Prove that there is a function $g$ such that
$f \circ g = I$.
\end{letters}

\pnum{*25.} Find a function $f$ such that $g \circ f = I$ for some $g$, but
such that there is no function $h$ with $f \circ h = I$.

\pnum{*26.} Suppose $f \circ g = I$ and $h \circ f = I$. Prove that $g=h$.
Hint: Use the fact that composition is associative.

\pnum{27.}
\begin{letters}
\item Suppose $f(x)=x+1$. Are there any functions $g$ such that
$f \circ g = g \circ f$?
\item Suppose $f$ is a constant function. For which functions $g$ does
$f \circ g = g \circ f$?
\item Suppose that $f \circ g = g \circ f$ for \emph{all} functions $g$. Show
that $f$ is the identity function, $f(x)=x$.
\end{letters}

\pnum{28.}
\begin{letters}
\item Let $F$ be the set of all functions whose domain is $\mathbf{R}$. Prove
that, using $+$ and $\cdot$ as defined in this chapter, all of properties
P1--P9 except P7 hold for $F$, provided $0$ and $1$ are interpreted as
constant functions.
\item Show that P7 does not hold.
\item[*(c)] Show that P10--P12 cannot hold. In other words, show that there
is no collection $P$ of functions in $F$, such that P10--P12 hold for $P$. (It
is sufficient, and will simplify things, to consider only functions which are
$0$ except at two points $x_0$ and $x_1$.)
\item[(d)] Suppose we define $f < g$ to mean that $f(x) < g(x)$ for all $x$.
Which of P$'$10--P$'$13 (in Problem~1-8) now hold?
\item[(e)] If $f < g$, is $h \circ f < h \circ g$? Is $f \circ h < g \circ h$?
\end{letters}

% Appendix. Ordered Pairs
\subsection*{APPENDIX. ORDERED PAIRS}

Not only in the definition of a function, but in other parts of the book as
well, it is necessary to use the notion of an ordered pair of objects. A
definition has not yet been given, and we have never even stated explicitly
what properties an ordered pair is supposed to have. The one property which
we will require states formally that the ordered pair $(a,b)$ should be
determined by $a$ and $b$, and the order in which they are given:
\[
\text{if }(a,b)=(c,d),\text{ then }a=c\text{ and }b=d.
\]

Ordered pairs may be treated most conveniently by simply introducing $(a,b)$
as an undefined term and adopting the basic property as an axiom---since this
property is the only significant fact about ordered pairs, there is not much
point worrying about what an ordered pair ``really'' is. Those who find this
treatment satisfactory need read no further.

The rest of this short appendix is for the benefit of those readers who will
feel uncomfortable unless ordered pairs are somehow defined so that this
basic property becomes a theorem. There is no point in restricting our
attention to ordered pairs of numbers; it is just as reasonable, and just as
important, to have available the notion of an ordered pair of any two
mathematical objects. This means that our definition ought to involve only
concepts common to all branches of mathematics. The one common concept which
pervades all areas of mathematics is that of a set, and ordered pairs (like
everything else in mathematics) can be defined in this context; an ordered
pair will turn out to be a set of a rather special sort.

The set $\{a,b\}$, containing the two elements $a$ and $b$, is an obvious
first choice, but will not do as a definition for $(a,b)$, because there is
no way of determining from $\{a,b\}$ which of $a$ or $b$ is meant to be the
first element. A more promising candidate is the rather startling set:
\[
\bigl\{\{a\},\ \{a,b\}\bigr\}.
\]
This set has two members, both of which are \emph{themselves} sets; one
member is the set $\{a\}$, containing the single member $a$, the other is the
set $\{a,b\}$. Shocking as it may seem, we are going to define $(a,b)$ to be
this set. The justification for this choice is given by the theorem
immediately following the definition---the definition works, and there really
isn't anything else worth saying.

\begin{spivakdefn}
\[
(a,b) = \bigl\{\{a\},\ \{a,b\}\bigr\}.
\]
\end{spivakdefn}

\begin{theorem}
If $(a,b)=(c,d)$, then $a=c$ and $b=d$.
\end{theorem}

\begin{proof}
The hypothesis means that
\[
\bigl\{\{a\},\ \{a,b\}\bigr\} = \bigl\{\{c\},\ \{c,d\}\bigr\}.
\]
Now $\{\{a\},\ \{a,b\}\}$ contains just two members, $\{a\}$ and $\{a,b\}$;
and $a$ is the only common element of these two members of
$\{\{a\},\ \{a,b\}\}$. Similarly, $c$ is the unique common member of both
members of $\{\{c\},\ \{c,d\}\}$. Therefore $a=c$. We therefore have
\[
\bigl\{\{a\},\ \{a,b\}\bigr\} = \bigl\{\{a\},\ \{a,d\}\bigr\},
\]
and only the proof that $b=d$ remains. It is convenient to distinguish $2$
cases.

\textit{Case 1.} $b=a$. In this case, $\{a,b\}=\{a\}$, so the set
$\{\{a\},\ \{a,b\}\}$ really has only one member, namely, $\{a\}$. The same
must be true of $\{\{a\},\ \{a,d\}\}$, so $\{a,d\}=\{a\}$, which implies that
$d=a=b$.

\textit{Case 2.} $b \neq a$. In this case, $b$ is in one member of
$\{\{a\},\ \{a,b\}\}$ but not in the other. It must therefore be true that $b$
is in one member of $\{\{a\},\ \{a,d\}\}$ but not in the other. This can
happen only if $b$ is in $\{a,d\}$, but $b$ is not in $\{a\}$; thus $b=a$ or
$b=d$, but $b \neq a$; so $b=d$.
\end{proof}
